% Figure: Lagrange-multiplier geometry. Level curves of f(x,y) and
% the constraint curve g(x,y)=0; at the constrained extremum the
% level curve is tangent to the constraint, i.e. grad f parallel to
% grad g.
\begin{figure}[htbp]
\centering
\begin{tikzpicture}[scale=1.0,
  fcurve/.style={engblue, thick},
  gcurve/.style={engred, very thick},
  gradf/.style={->, engblue, very thick},
  gradg/.style={->, engred, very thick},
  axis/.style={engblack, thick, ->},
  label/.style={font=\small},
]
% Axes
\draw[axis] (-0.3, 0) -- (4.3, 0) node[right, label] {$x$};
\draw[axis] (0, -0.3) -- (0, 3.3) node[above, label] {$y$};

% Level curves of f(x,y) = x + y = c, three values
\foreach \c in {1.0, 2.0, 2.8, 3.6} {
  \draw[fcurve, opacity=0.7] (\c, 0) -- (0, \c);
}
\node[label, engblue] at (3.5, 0.18) {$f = c$};

% Constraint curve g(x,y) = x^2 + y^2 - 4 = 0  (circle radius 2)
% Show the first-quadrant arc
\draw[gcurve, domain=0:90, samples=50] plot ({2*cos(\x)}, {2*sin(\x)});
\node[label, engred] at (1.6, 1.95) {$g = 0$};

% Optimum of f = x+y on the circle: x = y = sqrt(2), f = 2 sqrt(2) ~ 2.828
\coordinate (Pstar) at (1.4142, 1.4142);
\filldraw[engblack] (Pstar) circle (0.06);
\node[label, anchor=south west] at (Pstar) {$\vect{x}^{*}$};

% Gradient of f at the optimum: (1, 1), drawn scaled
\draw[gradf] (Pstar) -- ++(0.8, 0.8) node[anchor=south, label] {$\grad f$};
% Gradient of g at the optimum: (2x, 2y) = (2.83, 2.83), parallel to (1,1)
\draw[gradg] (Pstar) -- ++(1.1, 1.1) node[anchor=south west, label, xshift=4pt] {$\grad g$};

% A "wrong" point on the constraint where the gradients are not parallel
\coordinate (Q) at (1.732, 1.0);  % angle 30 degrees on the circle
\filldraw[engblack!50] (Q) circle (0.04);
\draw[gradf, opacity=0.6] (Q) -- ++(0.6, 0.6);
\draw[gradg, opacity=0.6] (Q) -- ++({1.732*0.4}, {1.0*0.4});
\node[label, anchor=west] at (Q) {\small not optimal};

\end{tikzpicture}
\caption[Lagrange-multiplier tangency.]{The constrained optimum of
$f$ on the curve $g = 0$ is the point where a level curve of $f$ is
tangent to the constraint (heavy black dot, with parallel gradients
$\grad f$ in blue and $\grad g$ in red). At any other point on the
constraint (grey, lower right), the two gradients are not parallel:
moving along the constraint in the direction that decreases the
angle between them changes $f$. The Lagrange-multiplier condition
$\grad f = \lambda \grad g$ encodes the parallelism algebraically,
with $\lambda$ the proportionality factor.}
\label{fig:vol02:ch11:lagrange-tangency}
\end{figure}
